\documentclass[varwidth=18cm, border=6pt]{standalone}
\usepackage{amsmath, amssymb}
\usepackage{tenkz}

\begin{document}

% T maps the open two-site MPDO tensor network to the open three-site
% MPDO tensor network, while S maps the three-site network back to the
% two-site network.  Each K has one virtual index on each horizontal side
% and one ket/bra pair vertically.  Each blue channel line lies between
% two complete tensor-network objects; it is not an additional contraction.
\[
  \begin{tenkzcd}[maps, species={channel}, column sep=10mm, row sep=8mm]
    \tnpic[inline, physical=updown]{
      \tn[box]{\mathcal K} & \tn[box]{\mathcal K}
    } &
    \tnpic[inline, physical=updown]{
      \tn[box]{\mathcal K} & \tn[box]{\mathcal K} & \tn[box]{\mathcal K}
    } \\
    \tnpic[inline, physical=updown]{
      \tn[box]{\mathcal K} & \tn[box]{\mathcal K} & \tn[box]{\mathcal K}
    } &
    \tnpic[inline, physical=updown]{
      \tn[box]{\mathcal K} & \tn[box]{\mathcal K}
    }
    \tnarrow[from={(1,1)}, to={(1,2)}, species=channel]{\mathcal T}
    \tnarrow[from={(2,1)}, to={(2,2)}, species=channel]{\mathcal S}
  \end{tenkzcd}
\]

% H^Omega_{k,h} is the middle-subspin Hilbert space.  The normalized active
% density and its completed extension act on this space and are adjoined by T_1.
\[
  \mathcal H^{\Omega}_{k,h}
  =\bigoplus_l(R_k\otimes L_l)\otimes(R_l\otimes L_h).
\]
On an active pair, let
\[
  \widehat\Omega_{k,h}
  =\frac{1}{a_kb_h}\bigoplus_l
    (\eta_{k,l}\otimes\eta_{l,h}),
  \qquad
  \overline\Omega_{k,h}=\widehat\Omega_{k,h}.
\]
On an inactive pair, $\overline\Omega_{k,h}$ denotes a chosen positive
trace-one density on $\mathcal H^{\Omega}_{k,h}$.  Thus, in every case,
\[
  \overline\Omega_{k,h}\in
    \operatorname{End}(\mathcal H^{\Omega}_{k,h}),
  \quad
  \operatorname{tr}(\overline\Omega_{k,h})=1.
\]

% T = T_2 T_1 T_0 maps the two-site sector coordinates R_2 to the
% three-site sector coordinates R_3.  Composition runs left-to-right.
% Every line has species=channel, so line ink records the common CPTP
% type; T_0, T_1, and T_2 name the constituent maps.
\[
  \begin{tenkzcd}[maps, species={channel}]
    \mathfrak R_2 &
    L_k\otimes R_h &
    (L_k\otimes R_h)\otimes\mathcal H^{\Omega}_{k,h} &
    \mathfrak R_3
    \tnarrow[from={(1,1)}, to={(1,2)}, species=channel]{\mathcal T_0}
    \tnarrow[from={(1,2)}, to={(1,3)}, species=channel]{\mathcal T_1}
    \tnarrow[from={(1,3)}, to={(1,4)}, species=channel]{\mathcal T_2}
  \end{tenkzcd}
\]

% Four fiberwise maps, one map per row.  Each row spells its complete
% domain and codomain; the parameters k and h remain symbolic.
\[
  \begin{tenkzcd}[maps, species={channel}]
    \mathfrak R_2 & L_k\otimes R_h \\
    (L_k\otimes R_h)\otimes(R_k\otimes L_h) & \mathfrak R_2 \\
    \mathfrak R_3 & L_k\otimes R_h \\
    (L_k\otimes R_h)\otimes\mathcal H^{\Omega}_{k,h} & \mathfrak R_3
    \tnarrow[from={(1,1)}, to={(1,2)}, species=channel]{\mathcal T_0}
    \tnarrow[from={(2,1)}, to={(2,2)}, species=channel]{\mathcal S_2}
    \tnarrow[from={(3,1)}, to={(3,2)}, species=channel]{\mathcal S_0}
    \tnarrow[from={(4,1)}, to={(4,2)}, species=channel]{\mathcal T_2}
  \end{tenkzcd}
\]

% T(R_2(K_2(X))) = R_3(K_3(X)); the channel line is the map T and the
% two endpoints are closed sector-coordinate matrices.
The refinement identity is
\(
  \begin{tenkzcd}[maps, inline, species={channel}]
    \mathfrak R_2({\cal K}_2(X)) & \mathfrak R_3({\cal K}_3(X))
    \tnarrow[from={(1,1)}, to={(1,2)}, species=channel]{\mathcal T}
  \end{tenkzcd}
\).

\end{document}
